% prelude.tex — shared macros for paper and presentation
% Usage: % prelude.tex — shared macros for paper and presentation
% Usage: % prelude.tex — shared macros for paper and presentation
% Usage: % prelude.tex — shared macros for paper and presentation
% Usage: \input{prelude} after \documentclass and package loading

% ── Glossary / acronyms ─────────────────────────────────────────────────────
\newacronym{sat}{SAT}{\emph{Boolean Satisfiability}}
\newacronym{maxSat}{maxSAT}{\emph{Maximum Satisfiability}}
\newacronym{pb}{PB}{\emph{Pseudo-Boolean}}
\newacronym{cp}{CP}{\emph{Constraint Programming}}
\newacronym{tsp}{TSP}{\emph{Travelling Salesperson Problem}}
\newacronym{mdd}{MDD}{\emph{Multivalued Decision Diagram}}
\newacronym{bdd}{BDD}{\emph{Binary Decision Diagram}}
\newacronym{csp}{CSP}{\emph{Constraint Satisfaction Problem}}
\newacronym{cop}{COP}{\emph{Constrained Optimisation Problem}}
\newacronym{mip}{MIP}{\emph{Mixed Integer Programming}}
\newacronym{ilp}{ILP}{\emph{Integer Linear Programming}}
\newacronym{lp}{LP}{\emph{Linear Programming}}
\newacronym{lbbd}{LBBD}{\emph{Logic-Based Benders Decomposition}}
\newacronym{mus}{MUS}{\emph{Minimal Unsatisfiable Subset}}

\newcommand{\milp}{\gls{ilp}\xspace}
\newcommand{\relaxation}{\acrshort{lp} relaxation\xspace}

% ── Text shorthands ──────────────────────────────────────────────────────────
\newcommand{\dquote}[1]{``#1''}
\newcommand{\quoted}[1]{`#1'}
\newcommand{\ie}{i.e.\xspace}
\newcommand{\eg}{e.g.\xspace}
\newcommand{\Eg}{E.g.\xspace}

% ── Author annotation commands (silenced by default) ─────────────────────────
\newcommand{\question}[1]{}
\newcommand{\pjs}[1]{}
\newcommand{\tias}[1]{}
\newcommand{\hb}[1]{}
\newcommand{\wout}[1]{}
\newcommand{\scrap}[1]{}
\newcommand{\ignore}[1]{}

\newcommand{\red}[1]{\textcolor{blue}{#1}}

% ── Math notation ────────────────────────────────────────────────────────────
\newcommand{\brackets}[1]{\ensuremath{[\![#1]\!]}}
\newcommand{\iverson}[1]{\ensuremath{\langle #1 \rangle}}
\newcommand{\beq}[2]{\brackets{#1 = #2}}
\newcommand{\bevarlookup}[1]{\ifnum9<1#1 \ifcase#1\or x\or y\or z\fi\else x_{#1}\fi}
\newcommand{\x}[1]{\bevarlookup{#1}}
\newcommand{\be}[2]{\beq{\bevarlookup{#1}}{#2}}
\newcommand{\besml}[2]{\scriptstyle\be{#1}{#2}}

\newcommand{\etAl}[1]{#1~\emph{et al.}}
\newcommand{\citeNoun}[2]{\etAl{#1}~\cite{#2}}

\DeclareMathOperator*{\argmax}{arg\,max}
\DeclareMathOperator*{\argmin}{arg\,min}

\newcommand{\assignment}{\ensuremath{\hat{v}}}
\newcommand{\true}{\mathit{true}}
\newcommand{\false}{\mathit{false}}
\newcommand{\Integers}{\mathbb{Z}}
\newcommand{\Reals}{\mathbb{R}}
\newcommand{\Bools}{\mathbb{B}}
\renewcommand{\emptyset}{\varnothing}

\newcommand{\vars}[1]{\ensuremath{vars(#1)}}
\newcommand{\scope}[1]{\ensuremath{scope(#1)}}
\newcommand{\cons}[1]{\ensuremath{cons(#1)}}
\newcommand{\bounds}[1]{\ensuremath{bounds(#1)}}

\newcommand{\interval}[2]{#1..#2}
\newcommand{\tuple}[1]{\langle #1 \rangle}
\newcommand{\sequence}[1]{[#1]}
\newcommand{\set}[1]{\{#1\}}
\newcommand{\setComp}[2]{\set{#1 ~|~ #2}}
\newcommand{\sequenceComp}[2]{\sequence{#1 ~|~ #2}}
\newcommand{\dom}[1]{D(#1)}
\newcommand{\pp}{+\!\!+}

\newcommand{\rhs}{s}

% ── Table-specific notation ──────────────────────────────────────────────────
\newcommand{\cols}{n}
\newcommand{\rows}{m}
\newcommand{\row}{r}
\newcommand{\col}{j}

\newcommand{\llinex}[1]{Line \ref{#1}}
\newcommand{\linex}[1]{line \ref{#1}}
\newcommand{\lines}[2]{lines \ref{#1}--\ref{#2}}

\newcommand{\ttable}{\texttt{table}\xspace}
\newcommand{\btable}{T^{\Bools}}
\newcommand{\btableTransposed}{\hat{T}^{\Bools}}
\newcommand{\transpose}[1]{\hat{#1}}

\newcommand{\nolbbd}[2]{#2}

% ── Configuration / approach labels ──────────────────────────────────────────
\newcommand{\textconfig}[1]{\textsc{#1}}
\newcommand{\baseGleb}{\textconfig{Int}\xspace}
\newcommand{\baseBool}{\textconfig{Bool}\xspace}

\newcommand{\baseMdd}{\textconfig{MDD}\xspace}
\newcommand{\baseMddInput}{\textconfig{\baseMdd-input}\xspace}
\newcommand{\baseMddDomIncr}{\textconfig{\baseMdd-dom}\xspace}
\newcommand{\baseMddGreedy}{\textconfig{\baseMdd-greedy}\xspace}

\newcommand{\lazyGenerate}{\textconfig{Generate}\xspace}
\newcommand{\lazyNone}{\textconfig{Generate}\xspace}
\newcommand{\lazyFractional}{\textconfig{Fractional}\xspace}
\newcommand{\lazyNegatives}{\textconfig{Negatives}\xspace}
\newcommand{\lazyCutlift}{\textconfig{Cutlift}\xspace}
\newcommand{\lazyShrink}{\textconfig{Shrink}\xspace}
\newcommand{\lazyCutoff}{\textconfig{Hybrid}\xspace}
\newcommand{\lazyCutoffT}[1]{\textconfig{Hybrid (#1)}\xspace}
\newcommand{\lazyAll}{\textconfig{All}\xspace}

% ── xy-pic node macros ───────────────────────────────────────────────────────
\newcommand{\xyo}[1]{*++[o][F-]{#1}}
\newcommand{\xyob}[1]{*++[o][F=]{#1}}
\newcommand{\xyoy}[1]{*++[o][F=]{#1}}
\newcommand{\xyoo}[1]{*++[o][F=]{#1}}

% ── MDD node labels (shared by \mddDiagram*, \mddFlowA--D) ────────────────────
\newcommand{\mddNodeRoot}{[]}
\newcommand{\mddNodeXtwo}{[2]}
\newcommand{\mddNodeXone}{[1]}
\newcommand{\mddNodeXfour}{[4]}
\newcommand{\mddNodeXthree}{[3]}
\newcommand{\mddNodeXtwoYone}{[2,1]}
\newcommand{\mddNodeMerge}{\scriptscriptstyle[1,2] \scriptscriptstyle[4,3]}
\newcommand{\mddNodeXthreeYtwo}{[3,2]}
\newcommand{\mddNodeSnk}{snk}

% ── Algorithm macros ─────────────────────────────────────────────────────────
\newcommand{\codeBlockDelim}{}
\newcommand{\tableCon}{\ensuremath{\texttt{table}(\hat{b}, \btable)}}
\newcommand{\generateInitArgs}{\ensuremath{\assignment, \tableCon}}
\newcommand{\generateArgs}{\ensuremath{\generateInitArgs, X, \rhs, R, V}}
\newcommand{\validCut}{\sum_{\col \in X} a_\col b_\col \leq \rhs}

% ── Cut lifting macros ──────────────────────────────────────────────────────
\newcommand{\Cols}{\interval{1}{\cols}}
\newcommand{\Rows}{\interval{1}{\rows}}
\newcommand{\colLift}{\col'}
\newcommand{\Rtight}{\ensuremath{R_{\text{tight}}}}
\newcommand{\Ntight}{\ensuremath{N_{\text{tight}}}}
\newcommand{\Xtight}{\ensuremath{X'}}
\newcommand{\cutliftThm}{%
Lifting is correct.
Suppose $\btable \models \sum_{\col \in X} a_\col b_\col \leq \rhs$  and $\btable \models b_{\colLift} \rightarrow \sum_{\col \in X} a_\col b_\col \leq \rhs-a_{\colLift}$ then
$\btable \models \sum_{\col \in X \cup \set{\colLift}} a_\col b_\col \leq \rhs$.%
}

% ── Negatives macros ─────────────────────────────────────────────────────────
\newcommand{\negate}[1]{\overline{#1}}
\newcommand{\partSize}[1]{\hat{P}_{#1}}
\newcommand{\unencoded}{B}

% ── Experimental evaluation ──────────────────────────────────────────────────
\newcommand{\timeout}{600\xspace}
\newcommand{\memout}{8\xspace}
\newcommand{\unk}{\texttt{UNK}\xspace}
\newcommand{\mem}{\texttt{MEM}\xspace}
\newcommand{\pst}{\texttt{POST}\xspace}
\newcommand{\sat}{\texttt{FEAS}\xspace}
\newcommand{\uns}{\texttt{UNS}\xspace}
\newcommand{\sol}{\texttt{SOLVED}\xspace}
\newcommand{\opt}{\texttt{OPT}\xspace}
\newcommand{\XCSPYears}{XCSP3'22\textendash25\xspace}
\newcommand{\thousands}{$10^3$}

% ── Figures ───────────────────────────────────────────────────────────────────
% Locally thicken xy-pic edges/frames within a group (restored at the closing brace)
\makeatletter
\newcommand{\xyThick}{\begingroup\xylinethick@=\dimexpr2\xylinethick@\relax}
\newcommand{\xyThickEnd}{\endgroup}
\makeatother
% MDD diagram for the running example (parameterised by xy spacing)
\newcommand{\mddDiagram}[2]{% #1 = @C (column sep), #2 = @R (row sep)
\xyThick
\xymatrix@C=#1@R=#2{
   \xyo{\mddNodeRoot}
       \ar[r]^{\be{1}{2}}
       \ar[rd]^{\be{1}{1}}
       \ar[rdd]^>>>>>>{\be{1}{4}}
       \ar[rddd]_{\be{1}{3}}
       & \xyo{\mddNodeXtwo} \ar[r]^{\be{2}{1}}
       & \xyo{\mddNodeXtwoYone} \ar[r]^{\be{3}{1}} \ar@/_2pc/[r]^{\be{3}{2}}
       & \xyoo{\mddNodeSnk} \\
   & \xyo{\mddNodeXone} \ar[rd]^{\be{2}{2}} &  \\
   & \xyo{\mddNodeXfour} \ar[r]_{\be{2}{3}}
       & \xyo{\mddNodeMerge} \ar[ruu]^{\be{3}{3}} \\
   & \xyo{\mddNodeXthree} \ar[r]^{\be{2}{2}}
       & \xyo{\mddNodeXthreeYtwo} \ar[ruuu]_{\be{3}{2}} \\
}%
\xyThickEnd
}
\newcommand{\mddDiagramHighlight}[2]{% same but with path 1,2,3 highlighted
\xyThick
\xymatrix@C=#1@R=#2{
   \xyo{\mddNodeRoot}
       \ar[r]^{\be{1}{2}}
       \ar@[blue][rd]^{\textcolor{blue}{\be{1}{1}}}
       \ar[rdd]^>>>>>>{\be{1}{4}}
       \ar[rddd]_{\be{1}{3}}
       & \xyo{\mddNodeXtwo} \ar[r]^{\be{2}{1}}
       & \xyo{\mddNodeXtwoYone} \ar[r]^{\be{3}{1}} \ar@/_2pc/[r]^{\be{3}{2}}
       & \xyoo{\mddNodeSnk} \\
   & \xyo{\textcolor{blue}{\mddNodeXone}} \ar@[blue][rd]^{\textcolor{blue}{\be{2}{2}}} &  \\
   & \xyo{\mddNodeXfour} \ar[r]_{\be{2}{3}}
       & \xyo{\textcolor{blue}{\mddNodeMerge}} \ar@[blue][ruu]^{\textcolor{blue}{\be{3}{3}}} \\
   & \xyo{\mddNodeXthree} \ar[r]^{\be{2}{2}}
       & \xyo{\mddNodeXthreeYtwo} \ar[ruuu]_{\be{3}{2}} \\
}%
\xyThickEnd
}
\newcommand{\mddDiagramHighlightBoth}[2]{% paths 1,2,3 (blue) and 4,3,3 (red)
\xyThick
\xymatrix@C=#1@R=#2{
   \xyo{\mddNodeRoot}
       \ar[r]^{\be{1}{2}}
       \ar@[blue][rd]^{\textcolor{blue}{\be{1}{1}}}
       \ar@[red][rdd]^>>>>>>{\textcolor{red}{\be{1}{4}}}
       \ar[rddd]_{\be{1}{3}}
       & \xyo{\mddNodeXtwo} \ar[r]^{\be{2}{1}}
       & \xyo{\mddNodeXtwoYone} \ar[r]^{\be{3}{1}} \ar@/_2pc/[r]^{\be{3}{2}}
       & \xyoo{\mddNodeSnk} \\
   & \xyo{\textcolor{blue}{\mddNodeXone}} \ar@[blue][rd]^{\textcolor{blue}{\be{2}{2}}} &  \\
   & \xyo{\textcolor{red}{\mddNodeXfour}} \ar@[red][r]_{\textcolor{red}{\be{2}{3}}}
       & \xyo{\textcolor{violet}{\mddNodeMerge}} \ar@[violet][ruu]^{\textcolor{violet}{\be{3}{3}}} \\
   & \xyo{\mddNodeXthree} \ar[r]^{\be{2}{2}}
       & \xyo{\mddNodeXthreeYtwo} \ar[ruuu]_{\be{3}{2}} \\
}%
\xyThickEnd
}

% ── Misc ─────────────────────────────────────────────────────────────────────
\newcommand{\omitVarJ}{[x_1, \ldots, x_{j-1}, x_{j+1}, \ldots, x_{n-1}, x_n]}
 after \documentclass and package loading

% ── Glossary / acronyms ─────────────────────────────────────────────────────
\newacronym{sat}{SAT}{\emph{Boolean Satisfiability}}
\newacronym{maxSat}{maxSAT}{\emph{Maximum Satisfiability}}
\newacronym{pb}{PB}{\emph{Pseudo-Boolean}}
\newacronym{cp}{CP}{\emph{Constraint Programming}}
\newacronym{tsp}{TSP}{\emph{Travelling Salesperson Problem}}
\newacronym{mdd}{MDD}{\emph{Multivalued Decision Diagram}}
\newacronym{bdd}{BDD}{\emph{Binary Decision Diagram}}
\newacronym{csp}{CSP}{\emph{Constraint Satisfaction Problem}}
\newacronym{cop}{COP}{\emph{Constrained Optimisation Problem}}
\newacronym{mip}{MIP}{\emph{Mixed Integer Programming}}
\newacronym{ilp}{ILP}{\emph{Integer Linear Programming}}
\newacronym{lp}{LP}{\emph{Linear Programming}}
\newacronym{lbbd}{LBBD}{\emph{Logic-Based Benders Decomposition}}
\newacronym{mus}{MUS}{\emph{Minimal Unsatisfiable Subset}}

\newcommand{\milp}{\gls{ilp}\xspace}
\newcommand{\relaxation}{\acrshort{lp} relaxation\xspace}

% ── Text shorthands ──────────────────────────────────────────────────────────
\newcommand{\dquote}[1]{``#1''}
\newcommand{\quoted}[1]{`#1'}
\newcommand{\ie}{i.e.\xspace}
\newcommand{\eg}{e.g.\xspace}
\newcommand{\Eg}{E.g.\xspace}

% ── Author annotation commands (silenced by default) ─────────────────────────
\newcommand{\question}[1]{}
\newcommand{\pjs}[1]{}
\newcommand{\tias}[1]{}
\newcommand{\hb}[1]{}
\newcommand{\wout}[1]{}
\newcommand{\scrap}[1]{}
\newcommand{\ignore}[1]{}

\newcommand{\red}[1]{\textcolor{blue}{#1}}

% ── Math notation ────────────────────────────────────────────────────────────
\newcommand{\brackets}[1]{\ensuremath{[\![#1]\!]}}
\newcommand{\iverson}[1]{\ensuremath{\langle #1 \rangle}}
\newcommand{\beq}[2]{\brackets{#1 = #2}}
\newcommand{\bevarlookup}[1]{\ifnum9<1#1 \ifcase#1\or x\or y\or z\fi\else x_{#1}\fi}
\newcommand{\x}[1]{\bevarlookup{#1}}
\newcommand{\be}[2]{\beq{\bevarlookup{#1}}{#2}}
\newcommand{\besml}[2]{\scriptstyle\be{#1}{#2}}

\newcommand{\etAl}[1]{#1~\emph{et al.}}
\newcommand{\citeNoun}[2]{\etAl{#1}~\cite{#2}}

\DeclareMathOperator*{\argmax}{arg\,max}
\DeclareMathOperator*{\argmin}{arg\,min}

\newcommand{\assignment}{\ensuremath{\hat{v}}}
\newcommand{\true}{\mathit{true}}
\newcommand{\false}{\mathit{false}}
\newcommand{\Integers}{\mathbb{Z}}
\newcommand{\Reals}{\mathbb{R}}
\newcommand{\Bools}{\mathbb{B}}
\renewcommand{\emptyset}{\varnothing}

\newcommand{\vars}[1]{\ensuremath{vars(#1)}}
\newcommand{\scope}[1]{\ensuremath{scope(#1)}}
\newcommand{\cons}[1]{\ensuremath{cons(#1)}}
\newcommand{\bounds}[1]{\ensuremath{bounds(#1)}}

\newcommand{\interval}[2]{#1..#2}
\newcommand{\tuple}[1]{\langle #1 \rangle}
\newcommand{\sequence}[1]{[#1]}
\newcommand{\set}[1]{\{#1\}}
\newcommand{\setComp}[2]{\set{#1 ~|~ #2}}
\newcommand{\sequenceComp}[2]{\sequence{#1 ~|~ #2}}
\newcommand{\dom}[1]{D(#1)}
\newcommand{\pp}{+\!\!+}

\newcommand{\rhs}{s}

% ── Table-specific notation ──────────────────────────────────────────────────
\newcommand{\cols}{n}
\newcommand{\rows}{m}
\newcommand{\row}{r}
\newcommand{\col}{j}

\newcommand{\llinex}[1]{Line \ref{#1}}
\newcommand{\linex}[1]{line \ref{#1}}
\newcommand{\lines}[2]{lines \ref{#1}--\ref{#2}}

\newcommand{\ttable}{\texttt{table}\xspace}
\newcommand{\btable}{T^{\Bools}}
\newcommand{\btableTransposed}{\hat{T}^{\Bools}}
\newcommand{\transpose}[1]{\hat{#1}}

\newcommand{\nolbbd}[2]{#2}

% ── Configuration / approach labels ──────────────────────────────────────────
\newcommand{\textconfig}[1]{\textsc{#1}}
\newcommand{\baseGleb}{\textconfig{Int}\xspace}
\newcommand{\baseBool}{\textconfig{Bool}\xspace}

\newcommand{\baseMdd}{\textconfig{MDD}\xspace}
\newcommand{\baseMddInput}{\textconfig{\baseMdd-input}\xspace}
\newcommand{\baseMddDomIncr}{\textconfig{\baseMdd-dom}\xspace}
\newcommand{\baseMddGreedy}{\textconfig{\baseMdd-greedy}\xspace}

\newcommand{\lazyGenerate}{\textconfig{Generate}\xspace}
\newcommand{\lazyNone}{\textconfig{Generate}\xspace}
\newcommand{\lazyFractional}{\textconfig{Fractional}\xspace}
\newcommand{\lazyNegatives}{\textconfig{Negatives}\xspace}
\newcommand{\lazyCutlift}{\textconfig{Cutlift}\xspace}
\newcommand{\lazyShrink}{\textconfig{Shrink}\xspace}
\newcommand{\lazyCutoff}{\textconfig{Hybrid}\xspace}
\newcommand{\lazyCutoffT}[1]{\textconfig{Hybrid (#1)}\xspace}
\newcommand{\lazyAll}{\textconfig{All}\xspace}

% ── xy-pic node macros ───────────────────────────────────────────────────────
\newcommand{\xyo}[1]{*++[o][F-]{#1}}
\newcommand{\xyob}[1]{*++[o][F=]{#1}}
\newcommand{\xyoy}[1]{*++[o][F=]{#1}}
\newcommand{\xyoo}[1]{*++[o][F=]{#1}}

% ── MDD node labels (shared by \mddDiagram*, \mddFlowA--D) ────────────────────
\newcommand{\mddNodeRoot}{[]}
\newcommand{\mddNodeXtwo}{[2]}
\newcommand{\mddNodeXone}{[1]}
\newcommand{\mddNodeXfour}{[4]}
\newcommand{\mddNodeXthree}{[3]}
\newcommand{\mddNodeXtwoYone}{[2,1]}
\newcommand{\mddNodeMerge}{\scriptscriptstyle[1,2] \scriptscriptstyle[4,3]}
\newcommand{\mddNodeXthreeYtwo}{[3,2]}
\newcommand{\mddNodeSnk}{snk}

% ── Algorithm macros ─────────────────────────────────────────────────────────
\newcommand{\codeBlockDelim}{}
\newcommand{\tableCon}{\ensuremath{\texttt{table}(\hat{b}, \btable)}}
\newcommand{\generateInitArgs}{\ensuremath{\assignment, \tableCon}}
\newcommand{\generateArgs}{\ensuremath{\generateInitArgs, X, \rhs, R, V}}
\newcommand{\validCut}{\sum_{\col \in X} a_\col b_\col \leq \rhs}

% ── Cut lifting macros ──────────────────────────────────────────────────────
\newcommand{\Cols}{\interval{1}{\cols}}
\newcommand{\Rows}{\interval{1}{\rows}}
\newcommand{\colLift}{\col'}
\newcommand{\Rtight}{\ensuremath{R_{\text{tight}}}}
\newcommand{\Ntight}{\ensuremath{N_{\text{tight}}}}
\newcommand{\Xtight}{\ensuremath{X'}}
\newcommand{\cutliftThm}{%
Lifting is correct.
Suppose $\btable \models \sum_{\col \in X} a_\col b_\col \leq \rhs$  and $\btable \models b_{\colLift} \rightarrow \sum_{\col \in X} a_\col b_\col \leq \rhs-a_{\colLift}$ then
$\btable \models \sum_{\col \in X \cup \set{\colLift}} a_\col b_\col \leq \rhs$.%
}

% ── Negatives macros ─────────────────────────────────────────────────────────
\newcommand{\negate}[1]{\overline{#1}}
\newcommand{\partSize}[1]{\hat{P}_{#1}}
\newcommand{\unencoded}{B}

% ── Experimental evaluation ──────────────────────────────────────────────────
\newcommand{\timeout}{600\xspace}
\newcommand{\memout}{8\xspace}
\newcommand{\unk}{\texttt{UNK}\xspace}
\newcommand{\mem}{\texttt{MEM}\xspace}
\newcommand{\pst}{\texttt{POST}\xspace}
\newcommand{\sat}{\texttt{FEAS}\xspace}
\newcommand{\uns}{\texttt{UNS}\xspace}
\newcommand{\sol}{\texttt{SOLVED}\xspace}
\newcommand{\opt}{\texttt{OPT}\xspace}
\newcommand{\XCSPYears}{XCSP3'22\textendash25\xspace}
\newcommand{\thousands}{$10^3$}

% ── Figures ───────────────────────────────────────────────────────────────────
% Locally thicken xy-pic edges/frames within a group (restored at the closing brace)
\makeatletter
\newcommand{\xyThick}{\begingroup\xylinethick@=\dimexpr2\xylinethick@\relax}
\newcommand{\xyThickEnd}{\endgroup}
\makeatother
% MDD diagram for the running example (parameterised by xy spacing)
\newcommand{\mddDiagram}[2]{% #1 = @C (column sep), #2 = @R (row sep)
\xyThick
\xymatrix@C=#1@R=#2{
   \xyo{\mddNodeRoot}
       \ar[r]^{\be{1}{2}}
       \ar[rd]^{\be{1}{1}}
       \ar[rdd]^>>>>>>{\be{1}{4}}
       \ar[rddd]_{\be{1}{3}}
       & \xyo{\mddNodeXtwo} \ar[r]^{\be{2}{1}}
       & \xyo{\mddNodeXtwoYone} \ar[r]^{\be{3}{1}} \ar@/_2pc/[r]^{\be{3}{2}}
       & \xyoo{\mddNodeSnk} \\
   & \xyo{\mddNodeXone} \ar[rd]^{\be{2}{2}} &  \\
   & \xyo{\mddNodeXfour} \ar[r]_{\be{2}{3}}
       & \xyo{\mddNodeMerge} \ar[ruu]^{\be{3}{3}} \\
   & \xyo{\mddNodeXthree} \ar[r]^{\be{2}{2}}
       & \xyo{\mddNodeXthreeYtwo} \ar[ruuu]_{\be{3}{2}} \\
}%
\xyThickEnd
}
\newcommand{\mddDiagramHighlight}[2]{% same but with path 1,2,3 highlighted
\xyThick
\xymatrix@C=#1@R=#2{
   \xyo{\mddNodeRoot}
       \ar[r]^{\be{1}{2}}
       \ar@[blue][rd]^{\textcolor{blue}{\be{1}{1}}}
       \ar[rdd]^>>>>>>{\be{1}{4}}
       \ar[rddd]_{\be{1}{3}}
       & \xyo{\mddNodeXtwo} \ar[r]^{\be{2}{1}}
       & \xyo{\mddNodeXtwoYone} \ar[r]^{\be{3}{1}} \ar@/_2pc/[r]^{\be{3}{2}}
       & \xyoo{\mddNodeSnk} \\
   & \xyo{\textcolor{blue}{\mddNodeXone}} \ar@[blue][rd]^{\textcolor{blue}{\be{2}{2}}} &  \\
   & \xyo{\mddNodeXfour} \ar[r]_{\be{2}{3}}
       & \xyo{\textcolor{blue}{\mddNodeMerge}} \ar@[blue][ruu]^{\textcolor{blue}{\be{3}{3}}} \\
   & \xyo{\mddNodeXthree} \ar[r]^{\be{2}{2}}
       & \xyo{\mddNodeXthreeYtwo} \ar[ruuu]_{\be{3}{2}} \\
}%
\xyThickEnd
}
\newcommand{\mddDiagramHighlightBoth}[2]{% paths 1,2,3 (blue) and 4,3,3 (red)
\xyThick
\xymatrix@C=#1@R=#2{
   \xyo{\mddNodeRoot}
       \ar[r]^{\be{1}{2}}
       \ar@[blue][rd]^{\textcolor{blue}{\be{1}{1}}}
       \ar@[red][rdd]^>>>>>>{\textcolor{red}{\be{1}{4}}}
       \ar[rddd]_{\be{1}{3}}
       & \xyo{\mddNodeXtwo} \ar[r]^{\be{2}{1}}
       & \xyo{\mddNodeXtwoYone} \ar[r]^{\be{3}{1}} \ar@/_2pc/[r]^{\be{3}{2}}
       & \xyoo{\mddNodeSnk} \\
   & \xyo{\textcolor{blue}{\mddNodeXone}} \ar@[blue][rd]^{\textcolor{blue}{\be{2}{2}}} &  \\
   & \xyo{\textcolor{red}{\mddNodeXfour}} \ar@[red][r]_{\textcolor{red}{\be{2}{3}}}
       & \xyo{\textcolor{violet}{\mddNodeMerge}} \ar@[violet][ruu]^{\textcolor{violet}{\be{3}{3}}} \\
   & \xyo{\mddNodeXthree} \ar[r]^{\be{2}{2}}
       & \xyo{\mddNodeXthreeYtwo} \ar[ruuu]_{\be{3}{2}} \\
}%
\xyThickEnd
}

% ── Misc ─────────────────────────────────────────────────────────────────────
\newcommand{\omitVarJ}{[x_1, \ldots, x_{j-1}, x_{j+1}, \ldots, x_{n-1}, x_n]}
 after \documentclass and package loading

% ── Glossary / acronyms ─────────────────────────────────────────────────────
\newacronym{sat}{SAT}{\emph{Boolean Satisfiability}}
\newacronym{maxSat}{maxSAT}{\emph{Maximum Satisfiability}}
\newacronym{pb}{PB}{\emph{Pseudo-Boolean}}
\newacronym{cp}{CP}{\emph{Constraint Programming}}
\newacronym{tsp}{TSP}{\emph{Travelling Salesperson Problem}}
\newacronym{mdd}{MDD}{\emph{Multivalued Decision Diagram}}
\newacronym{bdd}{BDD}{\emph{Binary Decision Diagram}}
\newacronym{csp}{CSP}{\emph{Constraint Satisfaction Problem}}
\newacronym{cop}{COP}{\emph{Constrained Optimisation Problem}}
\newacronym{mip}{MIP}{\emph{Mixed Integer Programming}}
\newacronym{ilp}{ILP}{\emph{Integer Linear Programming}}
\newacronym{lp}{LP}{\emph{Linear Programming}}
\newacronym{lbbd}{LBBD}{\emph{Logic-Based Benders Decomposition}}
\newacronym{mus}{MUS}{\emph{Minimal Unsatisfiable Subset}}

\newcommand{\milp}{\gls{ilp}\xspace}
\newcommand{\relaxation}{\acrshort{lp} relaxation\xspace}

% ── Text shorthands ──────────────────────────────────────────────────────────
\newcommand{\dquote}[1]{``#1''}
\newcommand{\quoted}[1]{`#1'}
\newcommand{\ie}{i.e.\xspace}
\newcommand{\eg}{e.g.\xspace}
\newcommand{\Eg}{E.g.\xspace}

% ── Author annotation commands (silenced by default) ─────────────────────────
\newcommand{\question}[1]{}
\newcommand{\pjs}[1]{}
\newcommand{\tias}[1]{}
\newcommand{\hb}[1]{}
\newcommand{\wout}[1]{}
\newcommand{\scrap}[1]{}
\newcommand{\ignore}[1]{}

\newcommand{\red}[1]{\textcolor{blue}{#1}}

% ── Math notation ────────────────────────────────────────────────────────────
\newcommand{\brackets}[1]{\ensuremath{[\![#1]\!]}}
\newcommand{\iverson}[1]{\ensuremath{\langle #1 \rangle}}
\newcommand{\beq}[2]{\brackets{#1 = #2}}
\newcommand{\bevarlookup}[1]{\ifnum9<1#1 \ifcase#1\or x\or y\or z\fi\else x_{#1}\fi}
\newcommand{\x}[1]{\bevarlookup{#1}}
\newcommand{\be}[2]{\beq{\bevarlookup{#1}}{#2}}
\newcommand{\besml}[2]{\scriptstyle\be{#1}{#2}}

\newcommand{\etAl}[1]{#1~\emph{et al.}}
\newcommand{\citeNoun}[2]{\etAl{#1}~\cite{#2}}

\DeclareMathOperator*{\argmax}{arg\,max}
\DeclareMathOperator*{\argmin}{arg\,min}

\newcommand{\assignment}{\ensuremath{\hat{v}}}
\newcommand{\true}{\mathit{true}}
\newcommand{\false}{\mathit{false}}
\newcommand{\Integers}{\mathbb{Z}}
\newcommand{\Reals}{\mathbb{R}}
\newcommand{\Bools}{\mathbb{B}}
\renewcommand{\emptyset}{\varnothing}

\newcommand{\vars}[1]{\ensuremath{vars(#1)}}
\newcommand{\scope}[1]{\ensuremath{scope(#1)}}
\newcommand{\cons}[1]{\ensuremath{cons(#1)}}
\newcommand{\bounds}[1]{\ensuremath{bounds(#1)}}

\newcommand{\interval}[2]{#1..#2}
\newcommand{\tuple}[1]{\langle #1 \rangle}
\newcommand{\sequence}[1]{[#1]}
\newcommand{\set}[1]{\{#1\}}
\newcommand{\setComp}[2]{\set{#1 ~|~ #2}}
\newcommand{\sequenceComp}[2]{\sequence{#1 ~|~ #2}}
\newcommand{\dom}[1]{D(#1)}
\newcommand{\pp}{+\!\!+}

\newcommand{\rhs}{s}

% ── Table-specific notation ──────────────────────────────────────────────────
\newcommand{\cols}{n}
\newcommand{\rows}{m}
\newcommand{\row}{r}
\newcommand{\col}{j}

\newcommand{\llinex}[1]{Line \ref{#1}}
\newcommand{\linex}[1]{line \ref{#1}}
\newcommand{\lines}[2]{lines \ref{#1}--\ref{#2}}

\newcommand{\ttable}{\texttt{table}\xspace}
\newcommand{\btable}{T^{\Bools}}
\newcommand{\btableTransposed}{\hat{T}^{\Bools}}
\newcommand{\transpose}[1]{\hat{#1}}

\newcommand{\nolbbd}[2]{#2}

% ── Configuration / approach labels ──────────────────────────────────────────
\newcommand{\textconfig}[1]{\textsc{#1}}
\newcommand{\baseGleb}{\textconfig{Int}\xspace}
\newcommand{\baseBool}{\textconfig{Bool}\xspace}

\newcommand{\baseMdd}{\textconfig{MDD}\xspace}
\newcommand{\baseMddInput}{\textconfig{\baseMdd-input}\xspace}
\newcommand{\baseMddDomIncr}{\textconfig{\baseMdd-dom}\xspace}
\newcommand{\baseMddGreedy}{\textconfig{\baseMdd-greedy}\xspace}

\newcommand{\lazyGenerate}{\textconfig{Generate}\xspace}
\newcommand{\lazyNone}{\textconfig{Generate}\xspace}
\newcommand{\lazyFractional}{\textconfig{Fractional}\xspace}
\newcommand{\lazyNegatives}{\textconfig{Negatives}\xspace}
\newcommand{\lazyCutlift}{\textconfig{Cutlift}\xspace}
\newcommand{\lazyShrink}{\textconfig{Shrink}\xspace}
\newcommand{\lazyCutoff}{\textconfig{Hybrid}\xspace}
\newcommand{\lazyCutoffT}[1]{\textconfig{Hybrid (#1)}\xspace}
\newcommand{\lazyAll}{\textconfig{All}\xspace}

% ── xy-pic node macros ───────────────────────────────────────────────────────
\newcommand{\xyo}[1]{*++[o][F-]{#1}}
\newcommand{\xyob}[1]{*++[o][F=]{#1}}
\newcommand{\xyoy}[1]{*++[o][F=]{#1}}
\newcommand{\xyoo}[1]{*++[o][F=]{#1}}

% ── MDD node labels (shared by \mddDiagram*, \mddFlowA--D) ────────────────────
\newcommand{\mddNodeRoot}{[]}
\newcommand{\mddNodeXtwo}{[2]}
\newcommand{\mddNodeXone}{[1]}
\newcommand{\mddNodeXfour}{[4]}
\newcommand{\mddNodeXthree}{[3]}
\newcommand{\mddNodeXtwoYone}{[2,1]}
\newcommand{\mddNodeMerge}{\scriptscriptstyle[1,2] \scriptscriptstyle[4,3]}
\newcommand{\mddNodeXthreeYtwo}{[3,2]}
\newcommand{\mddNodeSnk}{snk}

% ── Algorithm macros ─────────────────────────────────────────────────────────
\newcommand{\codeBlockDelim}{}
\newcommand{\tableCon}{\ensuremath{\texttt{table}(\hat{b}, \btable)}}
\newcommand{\generateInitArgs}{\ensuremath{\assignment, \tableCon}}
\newcommand{\generateArgs}{\ensuremath{\generateInitArgs, X, \rhs, R, V}}
\newcommand{\validCut}{\sum_{\col \in X} a_\col b_\col \leq \rhs}

% ── Cut lifting macros ──────────────────────────────────────────────────────
\newcommand{\Cols}{\interval{1}{\cols}}
\newcommand{\Rows}{\interval{1}{\rows}}
\newcommand{\colLift}{\col'}
\newcommand{\Rtight}{\ensuremath{R_{\text{tight}}}}
\newcommand{\Ntight}{\ensuremath{N_{\text{tight}}}}
\newcommand{\Xtight}{\ensuremath{X'}}
\newcommand{\cutliftThm}{%
Lifting is correct.
Suppose $\btable \models \sum_{\col \in X} a_\col b_\col \leq \rhs$  and $\btable \models b_{\colLift} \rightarrow \sum_{\col \in X} a_\col b_\col \leq \rhs-a_{\colLift}$ then
$\btable \models \sum_{\col \in X \cup \set{\colLift}} a_\col b_\col \leq \rhs$.%
}

% ── Negatives macros ─────────────────────────────────────────────────────────
\newcommand{\negate}[1]{\overline{#1}}
\newcommand{\partSize}[1]{\hat{P}_{#1}}
\newcommand{\unencoded}{B}

% ── Experimental evaluation ──────────────────────────────────────────────────
\newcommand{\timeout}{600\xspace}
\newcommand{\memout}{8\xspace}
\newcommand{\unk}{\texttt{UNK}\xspace}
\newcommand{\mem}{\texttt{MEM}\xspace}
\newcommand{\pst}{\texttt{POST}\xspace}
\newcommand{\sat}{\texttt{FEAS}\xspace}
\newcommand{\uns}{\texttt{UNS}\xspace}
\newcommand{\sol}{\texttt{SOLVED}\xspace}
\newcommand{\opt}{\texttt{OPT}\xspace}
\newcommand{\XCSPYears}{XCSP3'22\textendash25\xspace}
\newcommand{\thousands}{$10^3$}

% ── Figures ───────────────────────────────────────────────────────────────────
% Locally thicken xy-pic edges/frames within a group (restored at the closing brace)
\makeatletter
\newcommand{\xyThick}{\begingroup\xylinethick@=\dimexpr2\xylinethick@\relax}
\newcommand{\xyThickEnd}{\endgroup}
\makeatother
% MDD diagram for the running example (parameterised by xy spacing)
\newcommand{\mddDiagram}[2]{% #1 = @C (column sep), #2 = @R (row sep)
\xyThick
\xymatrix@C=#1@R=#2{
   \xyo{\mddNodeRoot}
       \ar[r]^{\be{1}{2}}
       \ar[rd]^{\be{1}{1}}
       \ar[rdd]^>>>>>>{\be{1}{4}}
       \ar[rddd]_{\be{1}{3}}
       & \xyo{\mddNodeXtwo} \ar[r]^{\be{2}{1}}
       & \xyo{\mddNodeXtwoYone} \ar[r]^{\be{3}{1}} \ar@/_2pc/[r]^{\be{3}{2}}
       & \xyoo{\mddNodeSnk} \\
   & \xyo{\mddNodeXone} \ar[rd]^{\be{2}{2}} &  \\
   & \xyo{\mddNodeXfour} \ar[r]_{\be{2}{3}}
       & \xyo{\mddNodeMerge} \ar[ruu]^{\be{3}{3}} \\
   & \xyo{\mddNodeXthree} \ar[r]^{\be{2}{2}}
       & \xyo{\mddNodeXthreeYtwo} \ar[ruuu]_{\be{3}{2}} \\
}%
\xyThickEnd
}
\newcommand{\mddDiagramHighlight}[2]{% same but with path 1,2,3 highlighted
\xyThick
\xymatrix@C=#1@R=#2{
   \xyo{\mddNodeRoot}
       \ar[r]^{\be{1}{2}}
       \ar@[blue][rd]^{\textcolor{blue}{\be{1}{1}}}
       \ar[rdd]^>>>>>>{\be{1}{4}}
       \ar[rddd]_{\be{1}{3}}
       & \xyo{\mddNodeXtwo} \ar[r]^{\be{2}{1}}
       & \xyo{\mddNodeXtwoYone} \ar[r]^{\be{3}{1}} \ar@/_2pc/[r]^{\be{3}{2}}
       & \xyoo{\mddNodeSnk} \\
   & \xyo{\textcolor{blue}{\mddNodeXone}} \ar@[blue][rd]^{\textcolor{blue}{\be{2}{2}}} &  \\
   & \xyo{\mddNodeXfour} \ar[r]_{\be{2}{3}}
       & \xyo{\textcolor{blue}{\mddNodeMerge}} \ar@[blue][ruu]^{\textcolor{blue}{\be{3}{3}}} \\
   & \xyo{\mddNodeXthree} \ar[r]^{\be{2}{2}}
       & \xyo{\mddNodeXthreeYtwo} \ar[ruuu]_{\be{3}{2}} \\
}%
\xyThickEnd
}
\newcommand{\mddDiagramHighlightBoth}[2]{% paths 1,2,3 (blue) and 4,3,3 (red)
\xyThick
\xymatrix@C=#1@R=#2{
   \xyo{\mddNodeRoot}
       \ar[r]^{\be{1}{2}}
       \ar@[blue][rd]^{\textcolor{blue}{\be{1}{1}}}
       \ar@[red][rdd]^>>>>>>{\textcolor{red}{\be{1}{4}}}
       \ar[rddd]_{\be{1}{3}}
       & \xyo{\mddNodeXtwo} \ar[r]^{\be{2}{1}}
       & \xyo{\mddNodeXtwoYone} \ar[r]^{\be{3}{1}} \ar@/_2pc/[r]^{\be{3}{2}}
       & \xyoo{\mddNodeSnk} \\
   & \xyo{\textcolor{blue}{\mddNodeXone}} \ar@[blue][rd]^{\textcolor{blue}{\be{2}{2}}} &  \\
   & \xyo{\textcolor{red}{\mddNodeXfour}} \ar@[red][r]_{\textcolor{red}{\be{2}{3}}}
       & \xyo{\textcolor{violet}{\mddNodeMerge}} \ar@[violet][ruu]^{\textcolor{violet}{\be{3}{3}}} \\
   & \xyo{\mddNodeXthree} \ar[r]^{\be{2}{2}}
       & \xyo{\mddNodeXthreeYtwo} \ar[ruuu]_{\be{3}{2}} \\
}%
\xyThickEnd
}

% ── Misc ─────────────────────────────────────────────────────────────────────
\newcommand{\omitVarJ}{[x_1, \ldots, x_{j-1}, x_{j+1}, \ldots, x_{n-1}, x_n]}
 after \documentclass and package loading

% ── Glossary / acronyms ─────────────────────────────────────────────────────
\newacronym{sat}{SAT}{\emph{Boolean Satisfiability}}
\newacronym{maxSat}{maxSAT}{\emph{Maximum Satisfiability}}
\newacronym{pb}{PB}{\emph{Pseudo-Boolean}}
\newacronym{cp}{CP}{\emph{Constraint Programming}}
\newacronym{tsp}{TSP}{\emph{Travelling Salesperson Problem}}
\newacronym{mdd}{MDD}{\emph{Multivalued Decision Diagram}}
\newacronym{bdd}{BDD}{\emph{Binary Decision Diagram}}
\newacronym{csp}{CSP}{\emph{Constraint Satisfaction Problem}}
\newacronym{cop}{COP}{\emph{Constrained Optimisation Problem}}
\newacronym{mip}{MIP}{\emph{Mixed Integer Programming}}
\newacronym{ilp}{ILP}{\emph{Integer Linear Programming}}
\newacronym{lp}{LP}{\emph{Linear Programming}}
\newacronym{lbbd}{LBBD}{\emph{Logic-Based Benders Decomposition}}
\newacronym{mus}{MUS}{\emph{Minimal Unsatisfiable Subset}}

\newcommand{\milp}{\gls{ilp}\xspace}
\newcommand{\relaxation}{\acrshort{lp} relaxation\xspace}

% ── Text shorthands ──────────────────────────────────────────────────────────
\newcommand{\dquote}[1]{``#1''}
\newcommand{\quoted}[1]{`#1'}
\newcommand{\ie}{i.e.\xspace}
\newcommand{\eg}{e.g.\xspace}
\newcommand{\Eg}{E.g.\xspace}

% ── Author annotation commands (silenced by default) ─────────────────────────
\newcommand{\question}[1]{}
\newcommand{\pjs}[1]{}
\newcommand{\tias}[1]{}
\newcommand{\hb}[1]{}
\newcommand{\wout}[1]{}
\newcommand{\scrap}[1]{}
\newcommand{\ignore}[1]{}

\newcommand{\red}[1]{\textcolor{blue}{#1}}

% ── Math notation ────────────────────────────────────────────────────────────
\newcommand{\brackets}[1]{\ensuremath{[\![#1]\!]}}
\newcommand{\iverson}[1]{\ensuremath{\langle #1 \rangle}}
\newcommand{\beq}[2]{\brackets{#1 = #2}}
\newcommand{\bevarlookup}[1]{\ifnum9<1#1 \ifcase#1\or x\or y\or z\fi\else x_{#1}\fi}
\newcommand{\x}[1]{\bevarlookup{#1}}
\newcommand{\be}[2]{\beq{\bevarlookup{#1}}{#2}}
\newcommand{\besml}[2]{\scriptstyle\be{#1}{#2}}

\newcommand{\etAl}[1]{#1~\emph{et al.}}
\newcommand{\citeNoun}[2]{\etAl{#1}~\cite{#2}}

\DeclareMathOperator*{\argmax}{arg\,max}
\DeclareMathOperator*{\argmin}{arg\,min}

\newcommand{\assignment}{\ensuremath{\hat{v}}}
\newcommand{\true}{\mathit{true}}
\newcommand{\false}{\mathit{false}}
\newcommand{\Integers}{\mathbb{Z}}
\newcommand{\Reals}{\mathbb{R}}
\newcommand{\Bools}{\mathbb{B}}
\renewcommand{\emptyset}{\varnothing}

\newcommand{\vars}[1]{\ensuremath{vars(#1)}}
\newcommand{\scope}[1]{\ensuremath{scope(#1)}}
\newcommand{\cons}[1]{\ensuremath{cons(#1)}}
\newcommand{\bounds}[1]{\ensuremath{bounds(#1)}}

\newcommand{\interval}[2]{#1..#2}
\newcommand{\tuple}[1]{\langle #1 \rangle}
\newcommand{\sequence}[1]{[#1]}
\newcommand{\set}[1]{\{#1\}}
\newcommand{\setComp}[2]{\set{#1 ~|~ #2}}
\newcommand{\sequenceComp}[2]{\sequence{#1 ~|~ #2}}
\newcommand{\dom}[1]{D(#1)}
\newcommand{\pp}{+\!\!+}

\newcommand{\rhs}{s}

% ── Table-specific notation ──────────────────────────────────────────────────
\newcommand{\cols}{n}
\newcommand{\rows}{m}
\newcommand{\row}{r}
\newcommand{\col}{j}

\newcommand{\llinex}[1]{Line \ref{#1}}
\newcommand{\linex}[1]{line \ref{#1}}
\newcommand{\lines}[2]{lines \ref{#1}--\ref{#2}}

\newcommand{\ttable}{\texttt{table}\xspace}
\newcommand{\btable}{T^{\Bools}}
\newcommand{\btableTransposed}{\hat{T}^{\Bools}}
\newcommand{\transpose}[1]{\hat{#1}}

\newcommand{\nolbbd}[2]{#2}

% ── Configuration / approach labels ──────────────────────────────────────────
\newcommand{\textconfig}[1]{\textsc{#1}}
\newcommand{\baseGleb}{\textconfig{Int}\xspace}
\newcommand{\baseBool}{\textconfig{Bool}\xspace}

\newcommand{\baseMdd}{\textconfig{MDD}\xspace}
\newcommand{\baseMddInput}{\textconfig{\baseMdd-input}\xspace}
\newcommand{\baseMddDomIncr}{\textconfig{\baseMdd-dom}\xspace}
\newcommand{\baseMddGreedy}{\textconfig{\baseMdd-greedy}\xspace}

\newcommand{\lazyGenerate}{\textconfig{Generate}\xspace}
\newcommand{\lazyNone}{\textconfig{Generate}\xspace}
\newcommand{\lazyFractional}{\textconfig{Fractional}\xspace}
\newcommand{\lazyNegatives}{\textconfig{Negatives}\xspace}
\newcommand{\lazyCutlift}{\textconfig{Cutlift}\xspace}
\newcommand{\lazyShrink}{\textconfig{Shrink}\xspace}
\newcommand{\lazyCutoff}{\textconfig{Hybrid}\xspace}
\newcommand{\lazyCutoffT}[1]{\textconfig{Hybrid (#1)}\xspace}
\newcommand{\lazyAll}{\textconfig{All}\xspace}

% ── xy-pic node macros ───────────────────────────────────────────────────────
\newcommand{\xyo}[1]{*++[o][F-]{#1}}
\newcommand{\xyob}[1]{*++[o][F=]{#1}}
\newcommand{\xyoy}[1]{*++[o][F=]{#1}}
\newcommand{\xyoo}[1]{*++[o][F=]{#1}}

% ── MDD node labels (shared by \mddDiagram*, \mddFlowA--D) ────────────────────
\newcommand{\mddNodeRoot}{[]}
\newcommand{\mddNodeXtwo}{[2]}
\newcommand{\mddNodeXone}{[1]}
\newcommand{\mddNodeXfour}{[4]}
\newcommand{\mddNodeXthree}{[3]}
\newcommand{\mddNodeXtwoYone}{[2,1]}
\newcommand{\mddNodeMerge}{\scriptscriptstyle[1,2] \scriptscriptstyle[4,3]}
\newcommand{\mddNodeXthreeYtwo}{[3,2]}
\newcommand{\mddNodeSnk}{snk}

% ── Algorithm macros ─────────────────────────────────────────────────────────
\newcommand{\codeBlockDelim}{}
\newcommand{\tableCon}{\ensuremath{\texttt{table}(\hat{b}, \btable)}}
\newcommand{\generateInitArgs}{\ensuremath{\assignment, \tableCon}}
\newcommand{\generateArgs}{\ensuremath{\generateInitArgs, X, \rhs, R, V}}
\newcommand{\validCut}{\sum_{\col \in X} a_\col b_\col \leq \rhs}

% ── Cut lifting macros ──────────────────────────────────────────────────────
\newcommand{\Cols}{\interval{1}{\cols}}
\newcommand{\Rows}{\interval{1}{\rows}}
\newcommand{\colLift}{\col'}
\newcommand{\Rtight}{\ensuremath{R_{\text{tight}}}}
\newcommand{\Ntight}{\ensuremath{N_{\text{tight}}}}
\newcommand{\Xtight}{\ensuremath{X'}}
\newcommand{\cutliftThm}{%
Lifting is correct.
Suppose $\btable \models \sum_{\col \in X} a_\col b_\col \leq \rhs$  and $\btable \models b_{\colLift} \rightarrow \sum_{\col \in X} a_\col b_\col \leq \rhs-a_{\colLift}$ then
$\btable \models \sum_{\col \in X \cup \set{\colLift}} a_\col b_\col \leq \rhs$.%
}

% ── Negatives macros ─────────────────────────────────────────────────────────
\newcommand{\negate}[1]{\overline{#1}}
\newcommand{\partSize}[1]{\hat{P}_{#1}}
\newcommand{\unencoded}{B}

% ── Experimental evaluation ──────────────────────────────────────────────────
\newcommand{\timeout}{600\xspace}
\newcommand{\memout}{8\xspace}
\newcommand{\unk}{\texttt{UNK}\xspace}
\newcommand{\mem}{\texttt{MEM}\xspace}
\newcommand{\pst}{\texttt{POST}\xspace}
\newcommand{\sat}{\texttt{FEAS}\xspace}
\newcommand{\uns}{\texttt{UNS}\xspace}
\newcommand{\sol}{\texttt{SOLVED}\xspace}
\newcommand{\opt}{\texttt{OPT}\xspace}
\newcommand{\XCSPYears}{XCSP3'22\textendash25\xspace}
\newcommand{\thousands}{$10^3$}

% ── Figures ───────────────────────────────────────────────────────────────────
% Locally thicken xy-pic edges/frames within a group (restored at the closing brace)
\makeatletter
\newcommand{\xyThick}{\begingroup\xylinethick@=\dimexpr2\xylinethick@\relax}
\newcommand{\xyThickEnd}{\endgroup}
\makeatother
% MDD diagram for the running example (parameterised by xy spacing)
\newcommand{\mddDiagram}[2]{% #1 = @C (column sep), #2 = @R (row sep)
\xyThick
\xymatrix@C=#1@R=#2{
   \xyo{\mddNodeRoot}
       \ar[r]^{\be{1}{2}}
       \ar[rd]^{\be{1}{1}}
       \ar[rdd]^>>>>>>{\be{1}{4}}
       \ar[rddd]_{\be{1}{3}}
       & \xyo{\mddNodeXtwo} \ar[r]^{\be{2}{1}}
       & \xyo{\mddNodeXtwoYone} \ar[r]^{\be{3}{1}} \ar@/_2pc/[r]^{\be{3}{2}}
       & \xyoo{\mddNodeSnk} \\
   & \xyo{\mddNodeXone} \ar[rd]^{\be{2}{2}} &  \\
   & \xyo{\mddNodeXfour} \ar[r]_{\be{2}{3}}
       & \xyo{\mddNodeMerge} \ar[ruu]^{\be{3}{3}} \\
   & \xyo{\mddNodeXthree} \ar[r]^{\be{2}{2}}
       & \xyo{\mddNodeXthreeYtwo} \ar[ruuu]_{\be{3}{2}} \\
}%
\xyThickEnd
}
\newcommand{\mddDiagramHighlight}[2]{% same but with path 1,2,3 highlighted
\xyThick
\xymatrix@C=#1@R=#2{
   \xyo{\mddNodeRoot}
       \ar[r]^{\be{1}{2}}
       \ar@[blue][rd]^{\textcolor{blue}{\be{1}{1}}}
       \ar[rdd]^>>>>>>{\be{1}{4}}
       \ar[rddd]_{\be{1}{3}}
       & \xyo{\mddNodeXtwo} \ar[r]^{\be{2}{1}}
       & \xyo{\mddNodeXtwoYone} \ar[r]^{\be{3}{1}} \ar@/_2pc/[r]^{\be{3}{2}}
       & \xyoo{\mddNodeSnk} \\
   & \xyo{\textcolor{blue}{\mddNodeXone}} \ar@[blue][rd]^{\textcolor{blue}{\be{2}{2}}} &  \\
   & \xyo{\mddNodeXfour} \ar[r]_{\be{2}{3}}
       & \xyo{\textcolor{blue}{\mddNodeMerge}} \ar@[blue][ruu]^{\textcolor{blue}{\be{3}{3}}} \\
   & \xyo{\mddNodeXthree} \ar[r]^{\be{2}{2}}
       & \xyo{\mddNodeXthreeYtwo} \ar[ruuu]_{\be{3}{2}} \\
}%
\xyThickEnd
}
\newcommand{\mddDiagramHighlightBoth}[2]{% paths 1,2,3 (blue) and 4,3,3 (red)
\xyThick
\xymatrix@C=#1@R=#2{
   \xyo{\mddNodeRoot}
       \ar[r]^{\be{1}{2}}
       \ar@[blue][rd]^{\textcolor{blue}{\be{1}{1}}}
       \ar@[red][rdd]^>>>>>>{\textcolor{red}{\be{1}{4}}}
       \ar[rddd]_{\be{1}{3}}
       & \xyo{\mddNodeXtwo} \ar[r]^{\be{2}{1}}
       & \xyo{\mddNodeXtwoYone} \ar[r]^{\be{3}{1}} \ar@/_2pc/[r]^{\be{3}{2}}
       & \xyoo{\mddNodeSnk} \\
   & \xyo{\textcolor{blue}{\mddNodeXone}} \ar@[blue][rd]^{\textcolor{blue}{\be{2}{2}}} &  \\
   & \xyo{\textcolor{red}{\mddNodeXfour}} \ar@[red][r]_{\textcolor{red}{\be{2}{3}}}
       & \xyo{\textcolor{violet}{\mddNodeMerge}} \ar@[violet][ruu]^{\textcolor{violet}{\be{3}{3}}} \\
   & \xyo{\mddNodeXthree} \ar[r]^{\be{2}{2}}
       & \xyo{\mddNodeXthreeYtwo} \ar[ruuu]_{\be{3}{2}} \\
}%
\xyThickEnd
}

% ── Misc ─────────────────────────────────────────────────────────────────────
\newcommand{\omitVarJ}{[x_1, \ldots, x_{j-1}, x_{j+1}, \ldots, x_{n-1}, x_n]}
